%% ma: L12-B08-0112
%% lop: 12
%% chuong: 2
%% bai: 8
%% dang: 12.8.8
%% loai: TN4
%% muc: TH
%% dap_an: "D"
%% nguon: "CĐ2.3.pdf giáo viên gửi 27/10/2026 (M07, Lê Minh Khôi – Nguyễn Anh Khôi, Chương II Chủ đề 3 Biểu thức tọa độ của các phép toán vectơ), Câu 24"
%% kien_thuc: "Bài 8 (điểm chia đoạn thẳng theo tỉ số, khoảng cách hai điểm)"
%% trang_thai: chinh-thuc
\begin{ex}
Trong không gian $Oxyz$, cho $A(2;0;0)$, $B(0;3;1)$, $C(-3;6;4)$. Gọi $M$ là điểm nằm trên đoạn thẳng $BC$ sao cho $MC=2MB$. Tính độ dài đoạn thẳng $AM$.
\choice{$AM=3\sqrt3$}{$AM=\sqrt{30}$}{$AM=2\sqrt7$}{\True $AM=\sqrt{29}$}
\loigiai{$\overrightarrow{BM}=\dfrac13\overrightarrow{BC}=(-1;1;1)\Rightarrow M(-1;4;2)$.

$\overrightarrow{AM}=(-3;4;2)$, $AM=\sqrt{29}$.}
\end{ex}
